% Nouveau EGA V (cote 25): archive cote 25, site page 142 (PDF page 143).
% English translation of the checked French annotated master.
% AI-translated by GPT-6 Astra (OpenAI, ChatGPT Pro mode).
\scanpage{25/142}
\mspage{\hand{3}}

prescheme x over k, \var{\xout{finally the last assertion follows from the first}}{}{typewriter strike-out}{ega25_25_142-v0001}
\var{\xout{preceding one and from the isomorphism $\underline{m}_x$ analogous $gr^2(P^2_{X/k})(x)$ Sy}}{}{typewriter strike-out}{ega25_25_142-v0002}

Observe that $gr^1(\underline{P}^1_{X/Y}) \hand{\simeq} \underline{\hand{\Omega}}^1_{X/Y}$; consequently the restriction
$d^1\hand{\phi}\,|V(\hand{\phi})$ can be interpreted as a section of $\underline{\hand{\Omega}}^1_{X/Y} \otimes \underline{O}_{V(\hand{\phi})}$,
which is simply the restriction of $d_{X/Y}\,\hand{\phi}$ to $V(\hand{\phi})$. We
can then consider the \uline{prescheme of zeros of this section}, which
we shall denote by $V(\hand{\phi})^{sing}$ and whose underlying set is precisely
the set of zeros of $\hand{\phi}$ singular relative to Y by Prop. 5.
(NB if \var{\xout{\unsure{$\hand{\phi}$}}}{}{typewriter strike-out}{ega25_25_142-v0003} $\hand{\mu}$ is a section of a locally free Module \var{\inserted[typed]{of finite type}}{of finite type}{insertion}{ega25_25_142-v0004} $\underline{E}$ on a
\var{\margin{$\|$}}{\placedmargin{$\|$}}{marginal note placed}{ega25_25_142-v0005}\var{\inserted[typed]{prescheme}}{prescheme}{insertion}{ega25_25_142-v0006} \var{\xout{ringed space}}{}{typewriter strike-out}{ega25_25_142-v0007} X, one defines the subprescheme of
zeros of $\hand{\mu}$ in an obvious way, for example as defined by the image Ideal of $\overset{\hand{\scriptscriptstyle\vee}}{\underline{E}} \hand{\longrightarrow} \underline{O}_X$
transpose of $\hand{\mu}$; when $\underline{E} = \underline{O}_X{}^n$ and $\hand{\mu} = \var{\xout{\unsure{(\hand{\mu}_i)}}}{}{typewriter strike-out}{ega25_25_142-v0008}\,(\hand{\mu}_1,..,\hand{\mu}_n)$, this
ideal is simply $\hand{\sum} \hand{\mu}_i\,\underline{O}_X$, which defines $V(\hand{\mu}_1,..,\hand{\mu}_n)$). Taking
now the restriction $d^2\hand{\phi}\,|V(\hand{\phi})^{sing}$, and observing that
$gr^2(P^2_{X/Y}) \hand{\simeq} Sym\var{\xout{m}}{}{typewriter strike-out}{ega25_25_142-v0009}^2(\underline{\hand{\Omega}}^1_{X/Y})$, one finds a canonical section $M(\hand{\phi})$ of
$Sym^2(\underline{\hand{\Omega}}^{\,1}_{X/\unsure{Y|X}}) \otimes \underline{O}_{V^{sing}}$. One verifies immediately that (taking
\var{\margin{$\|$}}{\placedmargin{$\|$}}{marginal note placed}{ega25_25_142-v0010}points of X with values in fields ...) this section is precisely
the one that determines the quadratic forms considered in Def. 1
(in the case \var{\xout{\illegible{1}}}{}{typewriter strike-out}{ega25_25_142-v0012} of an $X_k$ obtained from X/S by Spec(k)$\hand{\rightarrow}$ S). One
deduces from this a description of the set $V(\hand{\phi})^{sup\ sing}$ in terms of
this section, as follows: interpreting $M(\hand{\phi})$ as definin
homomorphism
\[
M(\hand{\phi})\hand{'} : \ \hand{\mathcal{T}}_{X/Y} \otimes_{\unsure{\underline{O}_X}} \unsure{\underline{O}_{V^{sing}}} \hand{\longrightarrow} \underline{\hand{\Omega}}^1_{X/Y} \otimes_{\underline{O}_X} \underline{O}_{V^{sing}} \quad ,
\]
one must take the set of points at which this homomorphism
is not an isomorphism. This shows in particular that $V(\hand{\phi})^{sup\ sing}$
is a closed set. One can specify \var{\deleted{it}}{}{deletion}{ega25_25_142-v0013} \var{\inserted{the latter}}{the latter}{insertion}{ega25_25_142-v0014}\var{\deleted{, by}}{}{deletion}{ega25_25_142-v0015} introducing
\[
D(\hand{\phi}) = \text{dét}\ M(\hand{\phi}) \ \hand{\in} \ \hand{\Gamma}\,(\underline{\hand{\Omega}}^d_{X/Y}{}^{\otimes 2} \otimes \underline{O}_{V^{sing}}\ ) \quad ,
\]

\mspage{\hand{142}}
\tnote{25-142-1}{The source opens with lowercase x after prescheme, retained here. The canceled Sy is an incomplete source fragment. The French letter cancellation eu within valeurs has no English effect; values is used without inventing a spelling error. The word translated as definin is clipped as définissan at the scan's right edge, and no missing ending or article is supplied. The source deletes the comma and en before introduisant; the resulting construction is retained. Within the literal formula, dét means determinant.}
